\chapter{Evaluation}
\label{cha:evaluation}

% Nachweis, dass die Ziele erreicht wurden – ehrlich, mit Zahlen, mit Grenzen.
% Struktur: Aufbau der Evaluation → Ergebnisse → Diskussion → Limitationen.

\section{Evaluationsdesign}
\label{sec:eval-design}

% Was wurde wie gemessen? Testdaten, Umgebung, Teilnehmende, Ablauf.
Die Evaluation erfolgte auf einem System mit \ldots{} Als Testdaten dienten \ldots

\section{Ergebnisse}
\label{sec:ergebnisse}

Tabelle~\ref{tab:ergebnisse} fasst die Messwerte zusammen.

\begin{table}[htbp]
  \centering
  \caption{Messergebnisse (Mittelwert aus zehn Durchläufen)}
  \label{tab:ergebnisse}
  \begin{tabular}{lrr}
    \toprule
    Szenario & Dauer in s & Standardabweichung \\
    \midrule
    Kleines Projekt (10 Seiten) & 2{,}1 & 0{,}2 \\
    Mittleres Projekt (40 Seiten) & 4{,}8 & 0{,}4 \\
    Großes Projekt (120 Seiten) & 11{,}3 & 0{,}9 \\
    \bottomrule
  \end{tabular}
\end{table}

\subsection{Erfüllung der Anforderungen}
% Jede Anforderung aus Kapitel 3 abhaken oder begründen, warum nicht.
Die Muss-Anforderungen FA1 und FA2 (Tabelle~\ref{tab:anforderungen-funktional})
sind vollständig erfüllt; FA4 wurde \ldots

\section{Diskussion}
\label{sec:diskussion}

% Interpretation: Was bedeuten die Ergebnisse? Stimmen sie mit Erwartungen und
% verwandten Arbeiten (Abschnitt~\ref{sec:verwandte-arbeiten}) überein?
Die Ergebnisse zeigen, dass \ldots{} Dies deckt sich mit \textcite{knuth1984literate},
der \ldots

\section{Limitationen}
\label{sec:limitationen}

% Grenzen der Aussagekraft: kleine Stichprobe, nur eine Plattform, …
Die Aussagekraft ist begrenzt durch \ldots
